\documentclass[a4paper,12pt]{report}

\usepackage[portuguese]{babel}
\usepackage[utf8]{inputenc}
\usepackage{booktabs}
\usepackage{pgfplots}
\usepackage{pgf-pie}
\usepackage{graphicx}
\usepackage{hyperref}
\usepackage{listings}
\usepackage{titlesec}
\usepackage{amsmath}
\usepackage{xcolor}
\usepackage{float}
\usepackage{tikz}
\usepackage[nottoc]{tocbibind}
\usepackage[backend=biber, style=alphabetic]{biblatex}
\usepackage{tikz}
\usetikzlibrary{arrows.meta, positioning}
\addbibresource{bibliografia.bib}

\pgfplotsset{compat=1.18}


\lstset{
    language=C,
    basicstyle=\ttfamily\small,
    frame=single,
    breaklines=true,
    postbreak=\mbox{\textcolor{red}{$\hookrightarrow$}\space},
    numbers=left,
    numberstyle=\tiny\color{gray},
    keywordstyle=\color{blue},
    stringstyle=\color{green!50!black},
    commentstyle=\color{gray},
    identifierstyle=\color{black},
    backgroundcolor=\color{white},
}

\lstdefinestyle{CStyle}{
    language=C,
    basicstyle=\ttfamily\small,
    frame=single,
    breaklines=true,
    postbreak=\mbox{\textcolor{red}{$\hookrightarrow$}\space},
    numbers=left,
    numberstyle=\tiny\color{gray},
    keywordstyle=\color{blue},
    stringstyle=\color{green!50!black},
    commentstyle=\color{gray},
    identifierstyle=\color{black},
    backgroundcolor=\color{white},
}

\title{\textbf{Projeto EDA - Relatório Etapa 1}}
\author{
    Licenciatura em Engenharia de Sistemas Informáticos \\ 
    Laboratórios de Informática 2024-25 \\ Leonel Freitas (31520)\\
    \href{https://github.com/Leonel-Freitas64/ProjetoEDA}{Repositorio GitHub}}
\date{\today}

\begin{document}

\maketitle
\tableofcontents
\newpage

\chapter{Introdução}
Este relatório apresenta o trabalho desenvolvido no âmbito do Projeto EDA, cujo objetivo principal foi criar uma aplicação em linguagem C para a gestão de antenas numa matriz, cálculo de efeitos nefastos entre antenas da mesma frequência e funcionalidades adicionais como gravar dados da aplicação em ficheiro binario para ser acessível mais tarde.

\section{Objetivos}
\begin{itemize}
    \item Representar antenas numa matriz com base em ficheiros de entrada.
    \item Detetar interferências (efeitos nefastos) entre antenas próximas com a mesma frequência.
    \item Poder adicionar e/ou remover antenas no programa.
    \item Imprimir a matriz com diferentes modos de visualização.
    \item Garantir a integridade dos dados e tratar sobreposições.
\end{itemize}

\chapter{Estrutura do Sistema}
O projeto foi organizado numa estrutura modular, com ficheiros separados para cada funcionalidade principal. A hierarquia do projeto é a seguinte:

\begin{verbatim}
ProjetoEDA/
+-- resources/
|   +--Estruturas.h
|   +--Funções.txt
+-- src/
|   +-- Antenas.c
|   +-- Nefasto.c
|   +-- Funções.c
|   +-- Binario.c
|   +-- main.c
+-- doxdoc/
|   +-- mainEDA.tex
|   +-- Doxyfile
+-- doc/
|   +-- main.pdf
\end{verbatim}

\newpage

\section{Descrição dos Ficheiros}
\begin{itemize}
    \item \texttt{Antenas.c}: Implementação das funções relacionadas à manipulação de antenas.
    \item \texttt{Nefasto.c}: Implementação das funções para cálculo de efeitos nefastos.
    \item \texttt{Funções.c}: Funções principais do programa, incluindo menus e manipulação de dados.
    \item \texttt{Binario.c}: Funções para leitura e escrita de ficheiros binários.
    \item \texttt{main.c}: Ficheiro principal que coordena a execução do programa.
    \item \texttt{mainEDA.tex}: Relatório da Etapa 1 do projeto.
    \item \texttt{Doxyfile}: Configuração do Doxygen para geração de documentação.
\end{itemize}

\chapter{Estruturas de Dados}
\section{Estrutura Antena}
\begin{lstlisting}[style=CStyle]
typedef struct Antena {
    char frequencia;
    int x, y;
    struct Antena *next;
} Antena;
\end{lstlisting}
Esta estrutura representa uma antena com a sua frequência e coordenadas na matriz. Está ligada a outras antenas através de uma lista ligada.

\section{Estrutura Nefasto}
\begin{lstlisting}[style=CStyle]
typedef struct Nefasto {
    int x1, y1, x2, y2;
    struct Antena* antena1;
    struct Antena* antena2; 
    struct Nefasto *next;
} Nefasto;
\end{lstlisting}
Cada efeito nefasto é registado com as duas posições onde ele ocorre, a duas casas de distância de cada antena envolvida e as respetivas antentas.
\section{Diagrama das Estruturas do Programa}

\begin{figure}[!ht]
    \centering
    \resizebox{1\textwidth}{!}{%
    \begin{tikzpicture}
    \tikzstyle{every node}=[font=\LARGE]
    \draw [ line width=1pt ] (11.75,11.5) rectangle (20,17);
    \node [font=\LARGE] at (15.5,16.25) {   Struct Antena};
    \node [font=\LARGE] at (15.75,15) {char frequencia};
    \node [font=\LARGE] at (15.75,14) {int x, y};
    \draw [line width=1pt, ->, >=Stealth] (20,12.25) .. controls (30.5,20.5) and (22.25,20.75) .. (16,17) ;
    \draw [line width=1pt, ->, >=Stealth] (11.75,3.75) .. controls (6,6.25) and (1.5,9) .. (11.75,14.75) ;
    \draw [ line width=1pt ] (11.75,6.25) rectangle (20.5,1.25);
    \node [font=\LARGE] at (16,5.5) {Struct Nefasto};
    \node [font=\LARGE] at (16,4.75) {int x1, x2, y1, y2};
    \node [font=\LARGE] at (16,3.5) {struct Antena* a1};
    \node [font=\LARGE] at (16,2.75) {struct Antena* a2};
    \node [font=\LARGE] at (16,2) {struct Nefasto* next};
    \node [font=\LARGE] at (16.25,12.25) {struct Antena* next};
    \draw [line width=1pt, ->, >=Stealth] (20.5,1.75) .. controls (31,10) and (22.75,10.25) .. (16.5,6.5) ;
    \end{tikzpicture}
    }%
    
    \label{fig:my_label}
    \end{figure}

\chapter{Funcionalidades do Programa}
\section{Leitura de Antenas de Ficheiro}
\begin{lstlisting}[style=CStyle]
Antena *lerAntenasDoFicheiro(const char *nomeFicheiro, int *linha, int *coluna);
\end{lstlisting}
Esta função lê uma matriz de um ficheiro e regista as posições de todas as antenas. Ignora células vazias e contabiliza corretamente linhas e colunas úteis.

\section{Cálculo de Efeitos Nefastos}
\begin{lstlisting}[style=CStyle]
Nefasto* calcularEfeitoNefasto(Antena *lista, int linha, int coluna);
\end{lstlisting}
Esta função percorre todas as antenas da mesma frequência e, se a distância for menor ou igual a 3, regista os efeitos nefastos a duas células de distância de cada antena, respeitando os limites da matriz.
\section{Impressão Visual da Matriz}
\begin{lstlisting}[style=CStyle]
void imprimirMatrizAntenas(Antena *lista, Nefasto *efeitos, int colunas, int linhas, int modo);
\end{lstlisting}
A matriz pode ser visualizada em três modos:
\begin{itemize}
    \item Modo 1: apenas antenas.
    \item Modo 2: apenas efeitos nefastos.
    \item Modo 3: ambos.
\end{itemize}

\section{Não Implementado : Redimensionamento e Deslocamento}
\begin{lstlisting}[style=CStyle]
Antena* redimensionarMatriz(int *linhas, int *colunas, Antena *lista);
\end{lstlisting}
Esta funcionalidade permite:
\begin{itemize}
    \item Aumentar ou diminuir dimensões da matriz.
    \item Remover antenas fora dos novos limites.
    \item Deslocar todas as antenas para a direita/baixo, se desejado.
\end{itemize}
Apesar de não implementada, a função foi desenhada e documentada para futura implementação.

\section{Adicionar Antena}
\begin{lstlisting}[style=CStyle]
void adicionarAntena(Antena **lista, char frequencia, int x, int y);
\end{lstlisting}
Esta função adiciona uma nova antena à lista ligada, verificando se as coordenadas são válidas e se já existe uma antena na mesma posição.

\section{Liberar Antena}
\begin{lstlisting}[style=CStyle]
void liberarAntena(Antena **lista, int x, int y);
\end{lstlisting}
Remove uma antena da lista ligada com base nas coordenadas fornecidas. Caso a antena não exista, a função retorna sem alterações.

\section{Salvar e Ler Ficheiro Binário}
\begin{lstlisting}[style=CStyle]
void salvarBinario(Antena *lista, const char *nomeFicheiro);
Antena* lerBinario(const char *nomeFicheiro);
\end{lstlisting}
Estas funções permitem salvar a lista de antenas em um ficheiro binário e carregá-la posteriormente. O formato binário garante eficiência e integridade dos dados.

\section{Menus Principal e Secundário}
\subsection{Menu Principal}
O menu principal apresenta opções para:
\begin{itemize}
    \item Adicionar, remover ou listar antenas.
    \item Calcular e listar efeitos nefastos.
    \item Salvar ou carregar dados de ficheiros.
    \item Alterar dimensões da matriz.
\end{itemize}

\subsection{Menu Secundário}
O menu secundário é exibido após a seleção de uma antena e permite:
\begin{itemize}
    \item Alterar a frequência da antena.
    \item Mover a antena para novas coordenadas.
    \item Remover a antena selecionada.
\end{itemize}

\section{Tratamento de Erros}
O programa inclui verificações para:
\begin{itemize}
    \item Antenas fora dos limites da matriz.
    \item Linhas vazias ou mal formatadas nos ficheiros de entrada.
    \item Coordenadas inválidas ao adicionar ou mover antenas.
    \end{itemize}
    Mensagens de erro claras são exibidas ao utilizador, garantindo uma experiência robusta.
    
    \chapter{Testes Realizados}
    Foram realizados vários testes com ficheiros contendo antenas em diferentes posições, com:
    \begin{itemize}
        \item Frequências iguais e distâncias menores ou maiores que 3.
        \item Efeitos diagonais, horizontais e verticais.
        \item Casos de sobreposição entre antena e efeito.
        \item Redimensionamentos e deslocamentos forçados.
    \end{itemize}
    
    \chapter{Conclusão}
    O projeto permitiu aplicar conhecimentos fundamentais de estruturas de dados, manipulação de ficheiros, listas ligadas e matrizes. A utilização de testes e visualização textual da matriz ajudou na validação dos resultados. O programa é modular, eficiente e extensível, permitindo futuras melhorias e adaptações.
    
    
    \end{document}